\section{Discussion}
\label{sec:discussion}

\paragraph{A protocol, not another motif count.}
The zero-common-relay counterexample identifies a distinction that motif
counts cannot express. Separate routes expose both sources in one
reconstruction, while the other loses a source at the readout. This
source-access difference is consistent with the archived high/chance
outcomes. It complements boundary-aware connectome controls~\cite{park2026}:
the routes used by a task can extend beyond the circuit family named by
an intervention. Interpreting a compensated deletion therefore requires
examining the replacement routes as well as the deleted structure.

\paragraph{Population information and finite trainability.}
Proposition~1 separates these quantities for the ideal task. All nonzero
$k=8$ dose reconstructions preserve the complete label information, yet
their finite-readout accuracies vary substantially. Their dose curves
therefore describe how the recorded learning procedure handles different
input mixtures within an information-preserving regime. Rank-one transfer
can also retain the label: full recovery of both source vectors is a
stronger requirement than this classification task needs. The useful
diagnostic is to test the task's information requirement before treating
a change in fitted accuracy as information deletion.

\paragraph{Scope.}
The archived task uses synthetic source labels, two-step linear propagation,
and a finite MLP on fixed subgraphs of one male connectome. The visual tests
use that propagator, the archived MLP, and the same three narrow
reconstructions. The four-source assay asks whether one CIFAR-10 class remains
readable when three other images drive the operator. It is not a
natural-image recognition benchmark. The two $k=8$
extractions share 195 neurons. The population result isolates source access
inside one reconstructed visual subgraph, rather than broader biological
organization. The exact population result assumes
noiseless continuous features; noise and finite precision affect subspace
decoding, while nonlinear or spiking propagation requires input-dependent
analysis. The separate reverse-support CIFAR-10 subset control detects no
connectome-specific accuracy benefit.

The two dense $k=8$ base graphs match frozen source hashes, but original
dose-instance hashes are unavailable. The reconstructed dose operators
are therefore post-hoc calculations, and their agreement with archived
outcomes is a consistency result. The edit algorithms also differ in order
and partner choices, preventing attribution of the historical score gap
to the common-relay rejection rule alone. Full provenance and the
recorded initialization protocol are given in the supplement.

\section{Conclusion}
\label{sec:conclusion}

Structural intervention reports should include the task information that survives.
For the fixed connectome propagator studied here, source visibility gives that
report, and the ideal subspace task turns it into an exact Bayes criterion.
On the matched narrow graphs, the archived MLP follows that matrix on the
subspace task and a visual parity XOR. With four simultaneous CIFAR-10 images, the same
matrix predicts which image class survives.
The same criterion separates matched zero-motif graphs, distinguishes ideal
information from finite-readout dose curves, and shows why a protected null
is a different counterfactual. The audit is specific to this operator and task,
but the reporting standard is broader: a preserved statistic is evidence only
for the constraint it actually enforces.
